\begin{abstract}

Large text embedding models are central to semantic search, clustering, and retrieval-augmented generation (RAG), where quality depends on the geometry induced by the teacher embedding space. Existing embedding distillation methods provide useful supervision through teacher vectors, token or layer representations, pairwise preferences, and mini-batch relations. These signals, however, do not explicitly describe the intrinsic manifold structure encoded by the teacher's neighborhoods: how local semantic regions connect, how similarity propagates through nearby examples, and which higher-order geometry should remain stable under compression. We propose \textbf{HeatGeo}, a heat-diffusion manifold distillation framework for compact text embedding models. HeatGeo caches teacher embeddings, constructs a mutual $k$-nearest-neighbor graph as a discrete semantic manifold, and derives multi-scale diffusion distributions that describe how semantic mass flows across this manifold. The student matches these diffusion distributions over sampled candidate neighborhoods, with a single pointwise anchor to cached teacher embeddings as the only auxiliary term. We show that a separate spectral objective is unnecessary: because the heat kernel is a spectral low-pass filter, matching diffusion at several scales already transfers the low-frequency topology that eigenvector regression would supply, and the resulting two-term objective needs no global eigendecomposition. HeatGeo therefore turns the teacher's final embeddings into geometry-aware supervision without requiring teacher hidden states, online teacher forward passes, or changes to the student architecture. We evaluate on STS, text classification, and pair classification to study whether manifold diffusion targets improve compact embedding distillation over pointwise, layer-level, and mini-batch relational baselines.

\end{abstract}
